\subsection{良序集的段}

一个关系 \(R\{x, y\}\) 称为 x 与 y 之间的良序关系，如果 R 是 x 与 y 之间的一个序关系，并且对于每个非空集合 E，若 \(R\{x, y\}\) 在 E 上诱导一个序关系（即满足 \(x \in E\) 蕴含 \(R\{x, x\}\) ；参见§1，第1号），则按此关系有序的 E 有一个最小元素。

一个由序 \(\Gamma\) 所排序的集合 E 被称为良序的，如果关系 \(y \in \Gamma\langle x \rangle\) 是 x 与 y 之间的良序关系； \(\Gamma\) 则称为 E 上的一个良序。以下定义与此等价：

定义1. 若集合E是有序的，且E的每个非空子集都有最小元素，则称E为良序的。

一个良序集 E 是全序的，因为 E 的每个子集 \(\{x, y\}\) 都有一个最小元素。E 的每个在 E 中有上界的子集 A 在 E 中都有一个最小上界。

\minorheading{示例}

\begin{enumerate}
\def\labelenumi{(\arabic{enumi})}
\item
  设 \(\mathbf{E} = \{\alpha, \beta\}\) 是一个元素互不相同的集合。容易验证， \(\{(\alpha, \alpha), (\beta, \beta), (\alpha, \beta)\}\) 这个 \(\mathbf{E} \times \mathbf{E}\) 的子集是 \(\mathbf{E}\) 上一个良序的图.
\item
  良序集的每个子集（特别是空子集）在诱导序下都是良序的。
\end{enumerate}

* (3) 存在非良序的全序集等价于无穷公理（§ 4，第 4 号，命题 3 的推论 1，以及练习 3）。

\begin{enumerate}
\def\labelenumi{(\arabic{enumi})}
\setcounter{enumi}{3}
\item
  如果 \(\Gamma\) 是 \(E\) 上的一个良序，则与 \(\Gamma\) 相反的序是 \(E\) 上的良序仅当 \(E\) 是有限的（ \(\S 4\) ，练习3）。*
\item
  设 E 是一个良序集。集合 \(E_{1}\) 通过向 E 添加一个最大元素 b（§1，第 7 号）而得到，是良序的，因为若 H 是 \(E_{1}\) 的任意非空子集且异于 \(\{b\}\) ，则 \(H \cap E\) 的最小元素也是 H 的最小元素。
\end{enumerate}

注记。 * 作为无穷公理（§6，第1号）的一个推论，存在没有最大元素的良序集，例如自然整数集 N。 *

定义2. 在有序集 E 中，E 的一个子集，若关系 \(x \in S\) , \(y \in E\) ，以及 \(y \leqslant x\) 蕴含 \(y \in S\) ，则称为E的一个段。

显然，E 的段的任意交和并都是 E 的一个段。若 S 是 E 的一个段，则 S 的每个段也是 E 的一个段。集合 E 本身和空集都是 E 的段。

命题1. 在一个良序集E中，E的每个不同于E本身的段都是一个区间 {]}←，a{[}，其中 a ∈ E。

设 S 为 E 的一个段，且 \(S \neq E\) 。由于 E --- S 非空，它有一个最小元素 a。根据定义 2，关系 \(x \geqslant a\) 蕴含 \(x \notin S\) ；否则就会有 \(a \in S\) ，这是矛盾的。因此 E --- S 是区间 \([a, \rightarrow[, \text{ 且 } S \text{ 是区间}] \leftarrow, a[\) .

对于每个元素 \(x\) （属于全序集 \(\mathbf{E}\) ），区间 \(] \leftarrow, x[\) 称为以 \(x\) 为端点的段，并记作 \(\mathbf{S}_x\) .

注意，若E是良序的且非空， \(S_{x}\) 有一个最小元素 \(\alpha\) ，因此它也是区间 \([\alpha, x[\) .

设E为全序集。诸 \(S_{x}\) （其中x遍历E）的并集A，当E无最大元时即为E；若E有最大元b，则 \(A = E - \{b\}\) .

命题2. 集合 \(\mathbf{E}^*\) （良序集 \(\mathbf{E}\) 的段所成的集合）按包含关系是良序的。映射 \(x \to S_x\) 是良序集 \(\mathbf{E}\) 到 \(\mathbf{E}\) 的异于 \(\mathbf{E}\) 本身的段所成集合上的同构。

显然，如果 \(x \in \mathbf{E}\) 并且 \(y \in \mathbf{E}\) ，关系 \(x \leqslant y\) 蕴含 \(\mathrm{S}_x \subset \mathrm{S}_y\) ，以及 \(x < y\) 蕴含 \(\mathrm{S}_x \neq \mathrm{S}_y\) ；因此映射 \(x \to \mathrm{S}_x\) 是 \(\mathbf{E}\) 到集合 \(\mathrm{S}(\mathbf{E})\) （ \(\mathbf{E}\) 的异于 \(\mathbf{E}\) 本身的段所成的集合）上的同构（§1，第12号，命题11），因此 \(\mathrm{S}(\mathbf{E})\) 是良序的。此外， \(\mathbf{E}^*\) 同构于由 \(\mathrm{S}(\mathbf{E})\) 通过添加一个最大元所得到的良序集。

命题3. 设 \((\mathbf{X}_{\iota})_{\iota \in \mathbf{I}}\) 是一族良序集，使得对于每一对指标 \((\iota, \varkappa)\) ，集合 \(\mathbf{X}_{\iota}, \mathbf{X}_{\varkappa}\) 中的一个是另一个的段。那么在集合 \(\mathbf{E} = \bigcup_{\iota \in \mathbf{I}} \mathbf{X}_{\iota}\) 上存在唯一的序关系，它在每一个 \(\mathbf{X}_{\iota}\) 上诱导出给定的序关系。在此序下， \(\mathbf{E}\) 是良序集。 \(X_{t}\) 的每一段是E的一个段；对于每个 \(x \in X_{t}\) ， \(X_{t}\) 中以x为端点的段等于E中以x为端点的段；并且E的每个段要么是E本身，要么是其中一个 \(X_{t}\) 的段。

第一个断言是以下一般引理的推论：

引理1. 设 \((\mathbf{X}_{\alpha})_{\alpha \in \Lambda}\) 为一个有序集族，它关于关系 \(c\) 是定向的（换言之，即对于每一对索引 \(\alpha, \beta\) 存在一个指标 \(\gamma\) 使得 \(\mathbf{X}_{\alpha} \subset \mathbf{X}_{\gamma}\) 并且 \(\mathbf{X}_{\beta} \subset \mathbf{X}_{\gamma}\) ）。假设对于每一对指标 \((\alpha, \beta)\) ，若 \(\mathbf{X}_{\alpha} \subset \mathbf{X}_{\beta}\) ，则 \(\mathbf{X}_{\alpha}\) 上由 \(\mathbf{X}_{\beta}\) 的序诱导的序与 \(\mathbf{X}_{\alpha}\) 上给定的序相同。在这些条件下，集合 \(\mathbf{E} = \bigcup_{\alpha \in \Lambda} \mathbf{X}_{\alpha}\) 上存在唯一的序关系，该序关系在每个子集 \(\mathbf{X}_{\alpha}\) 上诱导出给定的序关系。

设 \(G_{\alpha}\) 为给定序关系在 \(X_{\alpha}\) 上的图。若G是E上一个序的图，且该序在每一 \(X_{\alpha}\) 上诱导的序就是以 \(G_{\alpha}\) 为图的那个序，那么我们必须有 \(G_{\alpha} \subset G\) 对于每一个 \(\alpha \in A\) ；因此 G 包含 \(\bigcup_{\alpha \in A} G_{\alpha}\) 。另一方面，对于 E 的每一对元素 \((x, y)\) ，根据假设存在一个指标 \(\alpha \in A\) 使得 \(x \in X_{\alpha}\) 并且 \(y \in X_{\alpha}\) ；若 \((x, y) \in G\) ，我们有 \((x, y) \in G_{\alpha}\) ，因此 \(G \subset \bigcup_{\alpha \in A} G_{\alpha}\) 。因此，如果 E 上所需的序存在，则其图必然是 \(G = \bigcup_{\alpha \in A} G_{\alpha}\) 。剩下的只需证明该集合满足引理的条件。由于 \(G_{\beta} \cap (X_{\alpha} \times X_{\alpha}) = G_{\alpha}\) 若 \(X_{\alpha} \subset X_{\beta}\) ，我们有 \(G \cap (X_{\alpha} \times X_{\alpha}) = G_{\alpha}\) 对于所有 \(\alpha \in A\) ；另一方面，由假设可知，E 中任意三个元素 \(x, y, z\) 属于同一个 \(X_{\alpha}\) 。因此 \((x, y) \in G\) 是 E 上的一个序关系，引理得证。

我们现在开始证明命题3。首先，我们来证明每一个 \(\mathbf{X}_{\iota}\) 是 E 的一段。事实上，如果 \(x \in \mathbf{X}_{\iota}, y \in \mathbf{E}\) ，以及 \(y \leqslant x\) ，存在一个指标 \(x\) 使得 \(\mathbf{X}_{\iota} \subset \mathbf{X}_{x}\) 并且 \(y \in \mathbf{X}_{x}\) ；由于根据假设 \(\mathbf{X}_{\iota}\) 是 \(\mathbf{X}_{x}\) 的一段，我们有 \(y \in \mathbf{X}_{\iota}\) ，这证明了该断言。同样的推理证明，对于每一个 \(x \in \mathbf{X}_{\iota}\) ，以 \(x\) 为端点的段在 \(\mathbf{X}_{\iota}\) 中与E中的区间 {]}←, \(x[\) 相同。接下来让我们证明E是良序的。如果H是E的一个非空子集，则存在一个指标 \(\iota \in I\) 使得 \(\mathrm{H} \cap \mathrm{X}_{\iota} \neq \emptyset\) ；若 \(a\) 是 \(\mathrm{H} \cap \mathrm{X}_{\iota}\) 在 \(\mathbf{X}_{\iota}\) 中的最小元素，那么 \(a\) 也是 H 在 E 中的最小元素。也就是说，如果 \(x \in H\) ，存在一个指标 \(x \in I\) 使得 \(\mathbf{X}_{\iota} \subset \mathbf{X}_{x}\) 并且 \(x \in \mathbf{X}_{x}\) ；我们不可能有 \(x < a\) ，因为区间 {]}←, \(a[\) 包含于 \(\mathbf{X}_{\iota}\) ，因此我们有 \(x \geqslant a\) ，因为 \(\mathbf{X}_{x}\) 是全序的。

最后，我们必须证明E的一个异于E本身的段是其中一个 \(X_{t}\) 的段；这是前述论证的直接结果，因为这样的段具有形式 \(]\leftarrow\) , \(x[\text{(命题1)}\) ，并且由于x属于某个 \(X_{t}\) .

